% GENERATED FILE. Edit canonical lessons, then run npm run generate:books.
% canonical-source-hash: fnv1a64:15904a453df3254a
% canonical-lessons: IT-C465-lucertola, IT-C465-vipera, IT-C465-rospo, IT-C465-lumaca, IT-C465-strisciare

\chapter{Piccoli animali}
\label{ch:it-piccoli-animali}
\hlchaptermodality{465}

\begin{chapteropening}
\textbf{By the end of this chapter:} \emph{I can name small creatures on a walk: the lizard, the viper, the toad, the snail, and slithering.}

Complete the last lesson of chapter 465: I can name small creatures on a walk: the lizard, the viper, the toad, the snail, and slithering.
\end{chapteropening}

\section[la lucertola]{la lucertola — the lizard}

\label{lesson:IT-C465-lucertola}



\begin{quote}
Before the new one: say the Italian for the falcon, then the Italian for to graze.
\end{quote}

\subsection*{You'll want to know: la lucertola}
\textbf{la lucertola} — \textquotedblleft{}the lizard\textquotedblright{}.

\begin{quote}
Una lucertola prende il sole sul muro.
\end{quote}

\begin{quote}
\emph{A lizard is sunbathing on the wall.}
\end{quote}

Very common on old walls in summer.

\begin{cousinweb}
\emph{lucertola} is inherited from Latin \emph{lacertula}, a little \emph{lacerta} 'lizard' (EtymDB).
\end{cousinweb}

\begin{grammarlens}[title={What you've built}]
One more word for telling, explaining and describing at B1.
\end{grammarlens}

\subsection*{Your turn}
\begin{itemize}
\raggedright
  \item \emph{Say it:} \emph{la lucertola}
  \item \emph{Say it:} \emph{la lucertola}, once more
  \item \emph{Say it:} a sentence of your own with \emph{la lucertola}
  \item \emph{Recall:} say the Italian for the falcon, then the Italian for to graze, then say \emph{la lucertola} again
\end{itemize}

\begin{tcolorbox}[breakable,colback=teal!4,colframe=teal!35!black,title={Before you move on}]
What does la lucertola mean? (\textquotedblleft{}The lizard\textquotedblright{}.) Say it once more.
\end{tcolorbox}
\section[la vipera]{la vipera — the viper, the adder}

\label{lesson:IT-C465-vipera}



\begin{quote}
Before the new one: say the Italian for to graze, then the Italian for the lizard.
\end{quote}

\subsection*{You'll want to know: la vipera}
\textbf{la vipera} — \textquotedblleft{}the viper, the adder\textquotedblright{}.

\begin{quote}
In montagna attenzione alle vipere.
\end{quote}

\begin{quote}
\emph{In the mountains watch out for vipers.}
\end{quote}

Figuratively, a spiteful person.

\begin{cousinweb}
\emph{vipera} is a learned borrowing from Latin \emph{vīpera} 'viper' (EtymDB).
\end{cousinweb}

\begin{grammarlens}[title={What you've built}]
One more word for telling, explaining and describing at B1.
\end{grammarlens}

\subsection*{Your turn}
\begin{itemize}
\raggedright
  \item \emph{Say it:} \emph{la vipera}
  \item \emph{Say it:} \emph{la vipera}, once more
  \item \emph{Say it:} a sentence of your own with \emph{la vipera}
  \item \emph{Recall:} say the Italian for to graze, then the Italian for the lizard, then say \emph{la vipera} again
\end{itemize}

\begin{tcolorbox}[breakable,colback=teal!4,colframe=teal!35!black,title={Before you move on}]
What does la vipera mean? (\textquotedblleft{}The viper, the adder\textquotedblright{}.) Say it once more.
\end{tcolorbox}
\section[il rospo]{il rospo — the toad}

\label{lesson:IT-C465-rospo}



\begin{quote}
Before the new one: say the Italian for the lizard, then the Italian for the viper.
\end{quote}

\subsection*{You'll want to know: il rospo}
\textbf{il rospo} — \textquotedblleft{}the toad\textquotedblright{}.

\begin{quote}
Di sera i rospi cantano vicino allo stagno.
\end{quote}

\begin{quote}
\emph{In the evening the toads croak near the pond.}
\end{quote}

\emph{Ingoiare il rospo} — to swallow something unpleasant.

\begin{cousinweb}
The origin of \emph{rospo} is not settled.
\end{cousinweb}

\begin{grammarlens}[title={What you've built}]
One more word for telling, explaining and describing at B1.
\end{grammarlens}

\subsection*{Your turn}
\begin{itemize}
\raggedright
  \item \emph{Say it:} \emph{il rospo}
  \item \emph{Say it:} \emph{il rospo}, once more
  \item \emph{Say it:} a sentence of your own with \emph{il rospo}
  \item \emph{Recall:} say the Italian for the lizard, then the Italian for the viper, then say \emph{il rospo} again
\end{itemize}

\begin{tcolorbox}[breakable,colback=teal!4,colframe=teal!35!black,title={Before you move on}]
What does il rospo mean? (\textquotedblleft{}The toad\textquotedblright{}.) Say it once more.
\end{tcolorbox}
\section[la lumaca]{la lumaca — the snail, the slug}

\label{lesson:IT-C465-lumaca}



\begin{quote}
Before the new one: say the Italian for the viper, then the Italian for the toad.
\end{quote}

\subsection*{You'll want to know: la lumaca}
\textbf{la lumaca} — \textquotedblleft{}the snail, the slug\textquotedblright{}.

\begin{quote}
Dopo la pioggia escono le lumache.
\end{quote}

\begin{quote}
\emph{After the rain the snails come out.}
\end{quote}

\emph{Lento come una lumaca} — as slow as a snail.

\begin{cousinweb}
\emph{lumaca} is inherited from Vulgar Latin \emph{limāca}, from Latin \emph{līmāx} 'slug, snail' (EtymDB).
\end{cousinweb}

\begin{grammarlens}[title={What you've built}]
One more word for telling, explaining and describing at B1.
\end{grammarlens}

\subsection*{Your turn}
\begin{itemize}
\raggedright
  \item \emph{Say it:} \emph{la lumaca}
  \item \emph{Say it:} \emph{la lumaca}, once more
  \item \emph{Say it:} a sentence of your own with \emph{la lumaca}
  \item \emph{Recall:} say the Italian for the viper, then the Italian for the toad, then say \emph{la lumaca} again
\end{itemize}

\begin{tcolorbox}[breakable,colback=teal!4,colframe=teal!35!black,title={Before you move on}]
What does la lumaca mean? (\textquotedblleft{}The snail, the slug\textquotedblright{}.) Say it once more.
\end{tcolorbox}
\section[strisciare]{strisciare — to slither, to crawl; to swipe}

\label{lesson:IT-C465-strisciare}



\begin{quote}
Before the new one: say the Italian for the toad, then the Italian for the snail.
\end{quote}

\subsection*{You'll want to know: strisciare}
\textbf{strisciare} — \textquotedblleft{}to slither, to crawl; to swipe\textquotedblright{}.

\begin{quote}
Il serpente striscia tra i sassi.
\end{quote}

\begin{quote}
\emph{The snake slithers among the stones.}
\end{quote}

\emph{Strisciare la carta} — to swipe a card.

\begin{cousinweb}
\emph{strisciare} is built on \emph{striscia} 'stripe, trail', which EtymDB traces to Latin \emph{stria} 'furrow'; it is inherited.
\end{cousinweb}

\begin{grammarlens}[title={What you've built}]
One more word for telling, explaining and describing at B1.
\end{grammarlens}

\subsection*{Your turn}
\begin{itemize}
\raggedright
  \item \emph{Say it:} \emph{strisciare}
  \item \emph{Say it:} \emph{strisciare}, once more
  \item \emph{Say it:} a sentence of your own with \emph{strisciare}
  \item \emph{Recall:} say the Italian for the toad, then the Italian for the snail, then say \emph{strisciare} again
\end{itemize}

\begin{tcolorbox}[breakable,colback=teal!4,colframe=teal!35!black,title={Before you move on}]
What does strisciare mean? (\textquotedblleft{}To slither, to crawl; to swipe\textquotedblright{}.) Say it once more.
\end{tcolorbox}
